\section{Trabajos relacionados y fundamentos teóricos}
Las referencias lineales son necesarias al comparar modelos temporales profundos. DLinear y NLinear cuestionan que la complejidad arquitectónica explique por sí sola la precisión \cite{dlinear}. PatchTST usa parches independientes por canal \cite{patchtst}; TimesNet representa variación temporal en dos dimensiones \cite{timesnet}; iTransformer utiliza variables como dimensión de los elementos de atención \cite{itransformer}; TiDE emplea codificación densa \cite{tide}; y TimeMixer combina escalas temporales \cite{timemixer}. Estos trabajos tienen presupuestos y protocolos distintos. Aquí se ejecuta una configuración pequeña de PatchTST basada en código oficial; las otras cuatro arquitecturas son antecedentes bibliográficos, sin resultados experimentales propios en la cola ampliada.

Un dataset multivariante no obliga a que cada predictor utilice todas las variables conjuntamente. CPU y memoria son objetivos separados. La mayoría de modelos lineales comparte coeficientes entre servicios, pero consume la historia de un recurso de un servicio, normalizada con estadísticas propias. La GRU de grafo agrega vecinos. Por ello, una proyección univariante compartida no debe describirse como un modelo conjunto aprendido de toda la infraestructura. NLinear resta el último valor preparado, proyecta la historia centrada y vuelve a añadir ese valor; bajo ausencia, ese nivel reciente puede ser una imputación arrastrada, sin recuperación del dato oculto.

Una descomposición separa nivel suavizado y residual \cite{hyndmanfpp3}. El residual no demuestra estacionalidad, y una autocorrelación diaria tampoco basta para identificarla. MA25 repite extremos cuando opera sobre doce posiciones, de modo que el relleno forma parte del método. La componente exponencial actualiza un estado y la variante consciente de disponibilidad omite actualizaciones imputadas. Persistencia, media móvil y suavizamiento exponencial se evalúan también como predictores directos. Compartir una recurrencia no hace equivalentes un suavizador estadístico y una descomposición interna seguida de pesos entrenados.

Las GRU resumen secuencias mediante estados y compuertas \cite{gru}. GRU-D incorpora indicadores y decaimiento \cite{grud}; aquí se adapta a pesos compartidos entre servicios, sin reproducir su experimento clínico. BiTGraph y GinAR consideran estructuras temporales, relaciones y datos faltantes \cite{bitgraph,ginar}. Sus resultados históricos se preservan, incluso errores extremos, pero no se ajustan nuevamente en todos los escenarios. Las GRU local y de grafo amplían la comparación mediante información propia y agregados de llamadas. Una adyacencia estática conocida no demuestra influencia causal ni disponibilidad de cada arista en cada instante.

La imputación causal completa la entrada numérica; los indicadores registran si fue observada, arrastrada o reemplazada por mediana. La edad puede resumir una observación anterior al inicio de la ventana explícita. Rubin distingue los mecanismos de ausencia por la información de la que depende la disponibilidad \cite{rubinmissing}. El mecanismo sintético MAR dependiente del retardo vincula disponibilidad con CPU observada previamente, sin utilizar CPU contemporánea oculta. Su ejecución no prueba que las averías reales sigan MAR. La redundancia entre indicadores y edad también limita la interpretación aislada de coeficientes.

Los trabajos aplicados distinguen el objetivo predictivo y la acción posterior. TempoScale estudia predicción relacionada con autoescalado \cite{temposcale}, AERO combina predicción y orquestación \cite{aero}, y STAMP incorpora estado y dependencias espaciotemporales \cite{stamp}. Las revisiones de técnicas para microservicios y autoescalado sitúan estas contribuciones \cite{aimapping,autoscaling2025}. La RSL reciente de los autores aportó referencias candidatas; sus publicaciones primarias se verificaron antes de inclusión. El número de artículos de esa RSL no representa experimentos ejecutados aquí. Los trabajos sobre variables ausentes, muestreo dinámico, gestión adaptativa y anomalías aportan contexto adicional con su estado de publicación registrado \cite{ibn,radar2026,taxonomy2026,ibm2026}. Sus puntuaciones de anomalías no se comparan con MASE.

STEAM construye un grafo a partir de similitudes de carga y utiliza información del estado del sistema para obtener vistas alternativas del grafo \cite{steamgraph}. Este diseño permite contextualizar la adyacencia estática de llamadas que usamos aquí. DeepScaling evalúa una política de escalado orientada a estabilizar la utilización de CPU \cite{deepscaling}. Su evaluación incluye decisiones de control, mientras que este estudio mide el error de pronóstico sobre el uso registrado.

La incertidumbre de una diferencia media entre algoritmos es distinta de la incertidumbre de una futura observación. Los intervalos aquí remuestrean comparaciones bajo dependencia temporal \cite{blockbootstrap}; no son intervalos predictivos ni prueban cobertura. La predicción conforme requiere supuestos y calibraciones adicionales \cite{conformal}. MASE normaliza por variación del prefijo de entrenamiento \cite{mase}, pero no elimina diferencias semánticas entre CPU, memoria y utilización normalizada externa. La Tabla~\ref{tab:related} delimita el alcance implementado.
\begin{table*}[!t]
\caption{Estado del arte y alcance implementado. No se combinan resultados publicados en tareas distintas con las mediciones presentes.}\label{tab:related}
\centering\footnotesize
\begin{tabularx}{\linewidth}{p{.22\linewidth}XX}
\toprule Trabajo/familia & Mecanismo pertinente & Ejecución y límite \\\midrule
DLinear/NLinear \cite{dlinear} & Pronóstico afín y tratamiento del nivel & Ejecutados; historias cortas y selección fuente \\
GRU/GRU-D \cite{gru,grud} & Compuertas, disponibilidad y decaimiento & Adaptaciones ejecutadas; tarea clínica no reproducida \\
BiTGraph/GinAR \cite{bitgraph,ginar} & Grafos y variables incompletas & Solo históricos; presupuesto y tarea locales \\
PatchTST \cite{patchtst} & Parches temporales independientes por canal & Configuración pequeña ejecutada \\
TimesNet, iTransformer, TiDE, TimeMixer \cite{timesnet,itransformer,tide,timemixer} & Representaciones temporales alternativas & Solo revisados, sin clasificación experimental local \\
Predicción en nube \cite{powerprediction,temposcale,aero,stamp} & Integración con gestión de recursos & Contexto; no se midieron resultados de control \\
\bottomrule\end{tabularx}
\end{table*}
